	\documentclass{beamer}
\mode<presentation>
{
  \usetheme{myulm}
  \setbeamercovered{transparent}
  \setbeamertemplate{navigation symbols}{} % no navigation bar
  \setbeamersize{sidebar width left=1.17cm}
}

\usepackage[ngerman]{babel}
\usepackage[utf8]{inputenc}
\usepackage{amsmath,amssymb,amsfonts}
\usepackage{times}
\usepackage{graphicx}
\usepackage{fancyvrb}
\usepackage{array}
\usepackage{colortbl} % ING INF PSY
\usepackage{tikz}
\usepackage{listings}
\usepackage{verbatim}
\usetikzlibrary{arrows, decorations.markings}

% Anfang der Titelfolie
% Anpassung von: Titel, Untertitel, Autor, Datum und Institut

\title{Belief Propagation}
%\subtitle{Manchmal wird ein Untertitel ben\"{o}tigt, der ebenfalls einzeilig oder mehrzeilig sein kann}
\author{Michael Ruderer}
\newcommand{\presdatum}{\today} % alternativ zu \today: Eingabe eines festen Datums
\institute{Uni Ulm\\}
%Ende der Titelfolie


% Anfang der Kopfzeile der Folien
% Anpassung von: Zwischentitel, Leitthema oder Name
% Das Datum wird oben geändert: unter \presdatum{}!

\newcommand{\zwischentitel}{\insertsection}
\newcommand{\leitthema}{Belief Propagation}
% Ende der Kopfzeile

% Anfang der Folien
\begin{document}
\hspace*{-1.49cm}
\frame[plain]{\titlepage}

% Das Inhaltsverzeichnis
\hspace*{-0.7cm}
\section*{Inhalt/Überblick} % diese Section erscheint nicht im Inhaltsverzeichnis
\begin{frame}
  \frametitle{Table of Contents}
  \tableofcontents
\end{frame}


% 1. Folie
\section{Concepts}

\chapter{Belief Propagation}

\section{Factor Graphs}

For problems that include complex relations among many variables it is useful to have an abstract representation of how those	variables interact with each other. So called \emph{factor graphs} are such representations.\newline

In general, a factor graph represents the structure of a function's factorization into smaller functions. \newline If a function $f(X_1, \ldots, X_n)$ can be written as a product $\prod_{j=1}^{m}{f_j(S_j)}$  where the functions $f_j$ have smaller inputs $S_j \subset X$, its factorization 
can be expressed by a factor graph: The graph has two types of nodes: 
\emph{variable nodes} that correspond to the variables $X_i$ and \emph{factor nodes} corresponding to the functions $f_j$. An edge connects a variable node $X_i$ to a factor node $f_j$ if and only if $X_i$ is part of $f_j$'s input. 
This means the graph is an undirected bipartite graph with the node set $ V = \{X_1, \ldots, X_n\} \cup \{f_1, \ldots, f_n\}$ and edge set $E = \{(X_i, f_j) \; | \; X_i \in S_j\}$. 

In many applications the global function $f$ is a joined probability distribution that can be factorized by using information about independence between the variables. Typical tasks on factor graphs are computing variable assignments that maximize or minimize $f$ or computing marginal distributions if $f$ is a probability distribution. Both of these will be done in Section \ref{BPFS}.

\begin{example}
For random variables $X_1, X_2 \text{ and } X_3$ their joined probability distribution $f$ is defined as $f(x_1, x_2, x_3) := P(X_1 = x_1 \land X_2  = x_2 \land X_3 = x_3)$. 
$f$ can be factorized using the definition of conditional probabilities: $$f(x_1, x_2, x_3) = P(X_1 = x_1) * P(X_2 = x_2 \land X_3 = x_3 \; | \; X_1 = x_1)$$
This factorization still contains a function depending on all $3$ variables and is not useful. If $X_2$ and $X_3$ are known to be conditionally independent given $X_1$ this unwanted function can itself be factorized:

$$f(x_1, x_2, x_3) = \underbrace{P(X_1 = x_1)}_{f_1(x_1)} * \underbrace{P(X_2 = x_2 \; | \; X_1 = x_1)}_{f_2(x_1, x_2)} * \underbrace{P(X_3 = x_3 \; | \; X_1 = x_1)}_{f_3(x_1, x_3)}$$

In this factorization of $f$ each factor's input is smaller than the original one.

%\includegraphics[scale = 0.3]{img/Proba1}
\begin{figure}
\centering

\begin{tikzpicture}[scale=1.0,transform shape]
   	\node[shape=circle,draw=black] (x1) at (0,0) {$x_1$};
    \node[shape=circle,draw=black] (x2) at (-2,-2) {$x_2$};
    \node[shape=circle,draw=black] (x3) at (2,-2) {$x_3$};
    \node[rectangle,draw=black, label = {$f_1$}, fill] (f1) at (0,2) {};
    \node[rectangle, fill, draw=black, label = {$f_2$}] (f2) at (-1,-1) {};
    \node[rectangle, fill, draw=black, label = {$f_3$}] (f3) at (1, -1) {} ;

    \path [-] (x1) edge node[left] {} (f1);
    \path [-] (x1) edge node[left] {} (f2);
    \path [-] (x1) edge node[left] {} (f3);
    \path [-] (f2) edge node[left] {} (x2);
    \path [-] (f3) edge node[left] {} (x3);
   
\end{tikzpicture}
\caption{Factor graph of $f$'s factorization into $f_1, f_2, f_3$ \newline Usually variable nodes are drawn as circles whereas constraint nodes are drawn as rectangles to easily distinguish the types of nodes}

\end{figure}



\end{example}
% zu allgemein?
%
%A Constraint Satisfaction Problem (\textit{CSP}) consists of a set of variables $X = \{x_1, \ldots, x_n\}$, a set of domains $D = \{D_1, \ldots, D_n\}$ that specify which values each variable may take and a set of constraints $C = \{C_1, \ldots, C_m\}$ that must hold for the chosen values $x_i \in D_i$. \newline
%SAT can be written as CSP where each clause $C_j = {x_i, \ldots, x_k}$

%\newtheorem*{definition}{Definition}
\newpage
Factor graphs can also be used for describing constraint satisfaction problems. A factor corresponds to a constraint on its neighbour vertices, it evaluates to $1$ if the constraint is satisfied and to $0$ if not. For the global function $f$ - the product of all factors - to be $1$, every constraint has to be is satisfied. The special case of SAT problems is discussed in the following chapter.



\subsection{Factor graph of a SAT Problem}
Any SAT formula over the variables $x_1, \ldots, x_n$ in CNF form can be interpreted as a function that factorizes to the formulas clauses. \newline
$x_1, \ldots, x_n$ are boolean variables which can be either $0$ or $1$. A SAT formula $\mathcal{F}$ is a conjunction / product of special constraints, the formula's \emph{clauses}. A clause is a disjunction of variables and their negations  $(x_i\lor \overline{x_j} \lor \ldots)$ and is satisfied if at least one variable $x_i$ that appears unnegated is set to $1$ or at least one $x_j$ that appears negated is set to $0$. \newline
In the factor graph of a formula each factor node $a$ represents the local function defined by a single clause. If the variable $x_i$ or its negation $\overline{x_i}$ appears in this clause the factor graph contains an edge between $a$ and the variable node $i$. 

In following descriptions the letters $a$ and $b$ always describe factor nodes, $i$ and $j$ are used for variable nodes. Also often no distinction is made between the nodes of the factor graph and the corresponding elements of the formula. Especially from here on variables are denoted as single letters $i, j$ whereas the value of a variable $i$ is written as $x_i \in \{0, 1\}$

In \cite{survprob} some additional notation is defined to simplify the description of the algorithms in section \ref{BPFS}: 
\newpage
\begin{definition} Let $a$ be a factor node and $i$ one of its variables
\begin{itemize}
\item The value $J_i^a$ is $1$ if $i$ appears negated in $a$ and $-1$ if not
\item $V(a)$ is the set of all variables in $a$ \newline$V(i)$ is the set of all clauses containing $i$.  \newline Alternatively $V(a)$ and $V(i)$ are the neighbourhoods of $a$ and $i$ in the factor graph.
\item $V_+(a)$ is set of variables that appear unnegated in $a$ \newline $V_-(a)$ is set of $a$'s negated variables.
\item $V_a^u(i)$ is the set of $i$'s neighbours $b \neq a$ with $J_i^b \neq J_i^a$ \newline $V_a^s(i)$ is the set of neighbours $b \neq a$ with $J_i^b = J_i^a$
\end{itemize}
\end{definition}


\begin{example}
Let $$\mathcal{F} = \underbrace{(x_1 \lor x_2 \lor x_3)}_{a} \land \underbrace{(\overline{x_1} \lor x_2 \lor x_4) }_{b}\land \underbrace{(\overline{x_3} \lor x_4)}_{c} \land \underbrace{(\overline{x_1} \lor \overline{x_2})}_{d}$$

\begin{figure}[h]
\centering

\begin{tikzpicture}[scale=1.0,transform shape]
   	\node[shape=circle,draw=black] (x1) at (3,0) {$x_1$};
    \node[shape=circle,draw=black] (x2) at (1,-2) {$x_2$};
    \node[shape=circle,draw=black] (x3) at (-1,-2) {$x_3$};
    \node[shape=circle,draw=black] (x4) at (1,-4) {$x_4$};
    
    \node[rectangle,draw=black, label = {$a$}, fill] (a) at (0,0) {};
     \node[rectangle,draw=black, label = {$d$}, fill] (d) at (2,-1) {};
     \node[rectangle,draw=black, label = {$b$}, fill] (b) at (4,-2) {};
     \node[rectangle,draw=black, label = {$b$}, fill] (c) at (0,-3) {};
     
 
    \path [-] (x1) edge node[left] {} (a);
    \path [-] (x2) edge node[left] {} (a);
    \path [-] (x3) edge node[left] {} (a);
    
    \path [dashed] (x1) edge node[left] {} (d);
    \path [dashed] (x2) edge node[left] {} (d);
    
    \path [-] (x2) edge node[left] {} (b);
    \path [dashed] (x1) edge node[left] {} (b);
    \path [-] (x4) edge node[left] {} (b);
    
    \path [-] (x4) edge node[left] {} (c);
    \path [dashed] (x3) edge node[left] {} (c);
    

  
\end{tikzpicture}
\caption{Factor graph of a $\mathcal{F}$}
\end{figure}

Again the circles are variable nodes, the rectangles are constraint nodes. If if $x_i$ appears negated in the clause $a$ (if $J_a^i = 1$), the edge is drawn dashed.  

The variable sets of $b$ are for example $V(b) = \{1, 2, 4\}$, $V_+(b) = \{2, 4\}$, $ V_-(b) = \{1\}$.
The sets of $x_1$ are $V(1) = \{a, b, d\}$, $V_a^s(1) = \emptyset$, $V_a^u(1) = \{b, d\}$.

\end{example}



\section{Message Passing Algorithms on Trees}

If the factor graph is a tree, many problems can be solved efficiently using a form of dynamic programming called \emph{message passing}. \newline 
In general, message passing algorithms compute values for each edge of the factor graph. These values can be interpreted as \emph{messages} that are sent between the nodes. Since all edges connect factor nodes to variable nodes there can be two types of messages: messages passed from a variable $i$ to a factor $a$, denoted as $\mu_{i \rightarrow a}$ and messages passed from $a$ to $i$, denoted as $\mu_{a \rightarrow i}$. \newline
The messages must be defined so that a message $\mu_{a \rightarrow i}$ is determined by the messages $\mu_{j \rightarrow a}$ that $a$ received from neighbour variables $j \neq i$. 
The same must hold for $\mu_{i \rightarrow a}$. \newline
Usually the massages $\mu_{i \rightarrow a}$ are obtained by summing over $\mu_{b \rightarrow i}$ and $\mu_{a \rightarrow i}$ by multiplying the messages $\mu_{j \rightarrow a}$. The fundamental equation often has the form $$\mu_{a \rightarrow i} = \prod_{j \in V(a) \setminus i} \underbrace{\sum_{b \in V(j) \setminus a} \mu_{b \rightarrow j}}_{\mu_{j \rightarrow a}}$$ Therefore these types of algorithms are called \emph{sum-product}-algorithms. The reasoning behind the formula is explained in the description of the BP algorithm. The rule is formulated so that messages sent from factors to variables only appear as intermediate results in the computation of messages sent from variables to factors. \newline
For tree factor graphs which do not contain cycles the value of $\mu_{a \rightarrow i}$ does not influence its predecessors $\mu_{j \rightarrow a}$. The messages can therefore be computed sequentially starting with the factor graph's leaves.


\section{Message Passing on general graphs}
If a graph contains cycles, the above described messages are in general not well-defined. However, message passing algorithms which are correct for trees can be used as a heuristic for general graphs. \newline
This \emph{loopy} form of message passing consists of two parts. In the initialization  step each message is provisionally assigned a random value. Now that each message is set the sum-product equation can be used not to compute the absolute result but to repeatedly update the provisional values. In each update step the update rule is applied to all edges $a \rightarrow i$. There are several variants how to schedule these update tasks. For example the messages could be updated in parallel (\emph{synchronous}) or step by step (\emph{asynchronous}) in a fixed order. Here, the order of updates is chosen uniformly at random in each update step. If a message changed its value in the beginning of the step the following updates can instantly access and use its new value. 

The goal is to reach a point where no message would significantly change when applying the update rule to it. If this is the case, the messages are said to have \emph{converged}. A point of convergence is also called a \emph{fixed point}. Wether a fixed point is found  does not only depend on the input graph but also on the chosen starting values and the update schedule. Since there is no guarantee that the messages will converge it is possible for the update process to never terminate. In this case the computation has to be terminated after a fixed number of steps. Also an instance can have multiple fixed points that are not the desired solution. The converged messages are therefore only a heuristic solution if the factor graph contains cycles.


On trees this heuristic approach always converges to a unique set of fixed point warnings. The convergence for trees, like other properties of message passing algorithms, can be shown by induction on the \textit{level} of a message.

\begin{wrapfigure}{r}{0.3\textwidth}

\begin{tikzpicture}[scale=0.8,transform shape]
   	
	\node[shape=circle,draw=black] (u) at (0,0) {$u$};   	
	\node[shape=circle,draw=black] (v) at (2,0) {$v$};  
   	
  
	\node[] (x1) at (-1,0) {};
	\node[] (x2) at (-1,1) {};
	\node[] (x3) at (-1,-1) {};

    \draw[-] (u) edge [right] node {} (x1);
    \draw[-] (u) edge [right] node {} (x2);
    \draw[-] (u) edge [right] node {} (x3);
	
\tikzset{
    partial ellipse/.style args={#1:#2:#3}{
        insert path={+ (#1:#3) arc (#1:#2:#3)}
    }
}	
	
	\draw[-] (u) edge [right] node {} (v);
	\draw[dotted] (-0.5,0) [partial ellipse=-80:80:1cm and 1.5cm];
	\node[] (anon) at (0, 1.5) {$\tau_{u \rightarrow v}$};

\end{tikzpicture}
\end{wrapfigure}
\begin{definition} Let $(u, v)$ be an edge of a tree $\tau$. \newline
The tree $\tau_{u \rightarrow v}$ is the component of $\tau \setminus (u, v)$ that still contains $u$. The \emph{level} of the directed edge $u \rightarrow v$ is the height of $u$ in $\tau_{u \rightarrow v}$. 
\end{definition}
The level can be thought as the length of the longest path (the longest possible chain of messages) in $\tau$  that ends in $u$ and does not pass $v$.
\begin{lemma}\cite{survprob}
On a tree factor graph the message on an edge $a \rightarrow i$ at level $r$ converges after at most $r/2 + 1$ update steps.
\begin{proof} Induction on $r$:
\begin{itemize}
\item[] If $r = 0$, $a$ must be a leaf. Since there are no incoming messages $j \rightarrow a$ the message  $j \rightarrow i$ is constant at every step. \newline
Similarly for $r = 1$ all ingoing edges have level $0$ and $m_{a \rightarrow i}$ is constant from step $0$.
\item[] If $r \geq 2$ the message $m_{a \rightarrow i}$ is determined by the set of messages $m_{b \rightarrow j}$ sent to adjacent variables $j$. These messages are sent on edges with level $\leq r-2$ and have by induction converged on step $\frac{r-2}{2} + 1 = r / 2$. $m_{a \rightarrow i}$ converges in the next update step.
\end{itemize}
\end{proof}
\end{lemma}

\begin{figure}[h]
\centering

\begin{tikzpicture}[scale=0.8,transform shape]
   	\node[rectangle,draw=black, label = {$a$}, fill] (i) at (3,0) {$a$};
    \node[shape=circle,draw=black] (a) at (5,0) {$i$};
   
  \node[shape=circle,draw=black] (b1) at (0,2) {$j_1$};
  \node[shape=circle,draw=black] (b2) at (0,0) {$j_2$};
  \node[shape=circle,draw=black] (b3) at (0,-2) {$j_3$};

	\node[rectangle,draw=black, label = {$b_1$}, fill] (j1) at (-3,3) {};
	\node[rectangle,draw=black, label = {$b_2$}, fill] (j2) at (-3,1) {};
	\node[rectangle,draw=black, label = {$b_3$}, fill] (j3) at (-3,-1) {};
	\node[rectangle,draw=black, label = {$b_4$}, fill] (j4) at (-3,-3) {};


    
	\draw[->] (b3) edge [right] node {$m_{j_3 \rightarrow a}$} (i);
	\draw[->] (b2) edge [above] node {$m_{j_2 \rightarrow a}$} (i);
	\draw[->] (b1) edge [right] node {$m_{j_1 \rightarrow a}$} (i);
	\draw[->] (i) edge [above] node {$m_{a \rightarrow i}$} (a);
	
	\draw[->] (j1) edge [above] node {$m_{b_1 \rightarrow j_1}$} (b1);
  	\draw[->] (j2) edge [above] node {$m_{b_2 \rightarrow j_2}$} (b2);
	\draw[->] (j3) edge [above] node {$m_{b_3 \rightarrow j_2}$} (b2);
	\draw[->] (j4) edge [above] node {$m_{b_4 \rightarrow j_3}$} (b3);

\end{tikzpicture}
\caption{Subgraph of a factor tree with all messages required for computing $m_{a \rightarrow i}$}
\end{figure}

%\input{chapters(

\section{Warning Propagation}
\section{Warning Propagation}

The messages used in the \emph{Warning Propagation} Algorithm (WP) presented in \cite{survprob} are called \emph{warnings}. A warning $u_{a\rightarrow i} \in \{0, 1\}$ is passed from clause $a$ to variable $i$. A converged warning  $u^{\star}_{a\rightarrow i}$ with value $1$ should indicate, that  to satisfy the clause $a$, the variable $i$ has to take the value $1$ if $j \in V_+(a)$ or $0$ if $j \in V_-(a)$. Equivalently if $i$ does not satisfy $a$, the entire formula can not be satisfied. The warning $u^{\star}_{a \rightarrow i}$ will \emph{fix} the variable $i$. 
\subsection{Propagation Algorithm}

Like the general algorithm described in section \ref{BPA} warning propagation is correct on trees and can be used as a heuristic for cyclic graphs by randomly initializing the warnings and hoping for convergence.

The algorithm starts by assigning each warning $u_{a\rightarrow i}$ a random starting value and updates these provisional warnings until their values have converged to a set of fixed point warnings $u^{\star}_{a \rightarrow i}$ or until the number of iterations has exceeded some limit $t_{max}$.

The general idea is that the clause $a$ has to fix the variable $i$ only if all of its other variables $j \in V(a)\setminus i$ are already fixed to values that do not satisfy the clause $a$. \newline
The first step in the update procedure is to compute for each $j \in V(a) \setminus i$ the so called \emph{cavity field} $h_{j \rightarrow a}$ that indicates what value $j$ should take in the subproblem defined by $\tau_{j \rightarrow a}$. To compute $h_{j \rightarrow a}$ one has to count how many of the clauses $b \neq a$ fix $j$ to $1$ and how many fix $j$ to $0$:
$$h_{j \rightarrow a} = \sum_{b \in V_+(j)\setminus a}{u_{b \rightarrow j}} - \sum_{b \in V_-(j)\setminus a}{u_{b \rightarrow j}} = -\sum_{b \in V(j)} J^b_j u_{b \rightarrow j}$$

The clauses $b \in V_+(j)$ are the ones that would fix $j$ to $1$ if their warnings are active, the clauses $b \in V_-(j)$ would fix $j$ to $0$.
So if $h_{j \rightarrow a}$ is positive the variable $j$ tends to the value $1$. If the cavity field is negative it tends to $0$. If $h_{j \rightarrow a} = 0$ which includes the case $V(j) \setminus a = \emptyset$ no conclusion can be made.

When all cavity fields are computed, each variable $j$ with $h_{j \rightarrow a} \neq 0$ has a preferred value. This preferred value either makes the clause $a$ satisfied or does not contribute to its satisfaction. If all variables $j \in V(a)\setminus i$ prefer a non satisfying value, the clause $a$ sends a warning to $i$, meaning that $i$ should take the satisfying value.

This warning can be computed by $$u_{a \rightarrow i} = \prod_{j\in V(a)\setminus i}{\theta(h_{j \rightarrow a}J_j^a)} \quad \text{ with } \theta(x) = \left\{
  \begin{array}{@{}ll@{}}
    0, & \text{if}\ x \leq 0 \\
    1, & \text{otherwise}
  \end{array}\right\}$$

The factor $\theta(h_{j \rightarrow a}J_j^a)$ is $1$ if $j$ prefers to violate $a$ and $0$ if not:
\begin{itemize}
\item[] If $j$ has no preferred value its cavity field is $0$ and $\theta(h_{j \rightarrow a}J_j^a) = 0$ meaning no warning will be sent.
\item[] If the preferred value of $j$ satisfies $a$, $J_j^a$ and $h_{j \rightarrow a}$ have different signs and \newline $\theta(h_{j \rightarrow a}J_j^a) = \theta(-1)$ is again $0$.
\item[] If the preferred value of $j$ \emph{violates} $a$, $J_j^a$ and $h_{j \rightarrow a}$ have the same sign and $\theta(h_{j \rightarrow a}J_j^a) = 1$. If this is the case for all $j \neq i$ the product evaluates to $1$ and $a$ sends a warning to $i$.
\end{itemize}
\begin{lstlisting}[mathescape=true, frame = single]
	Warning Propagation Algorithm
	  Input: A CNF formula and its factor graph
	
	0. Randomly initialize all warnings $u_{a \rightarrow i} \in_R \{0, 1\}$
	
	1. For $t=0$ to $t = t_{max}$
		1.1 Compute in random order for all edges (a, i) $u_{a \rightarrow i} := \prod_{j \in V(a) \setminus i} {\theta \left( -\sum_{b\in V(j) \setminus a}{(J_j^a J_j^b) u_{b \rightarrow j}}  \right)}$
		1.2 If no message has changed goto 2.
	2. If $t = t_{max}$ return UN-CONVERGED, else return the generated warnings $u_{a \rightarrow i}^{\star}$
\end{lstlisting}
% $u_{a \rightarrow i} := \prod_{j \in V(a) \setminus i} {\theta \left( -J_j^a \left( \sum_{b\in V(j) \setminus a}{J_j^b u_{b \rightarrow j}} \right) \right)}$

The following lemma shows that - on trees - the computed messages indeed serve the purpose described in the first paragraph: If the algorithm successfully returns a set of converged warnings, each warning $u^{\star}_{a \rightarrow i} = 1$ fixes the variable $i$ to the value satisfying a:

\begin{lemma}\cite{survprob} Let $a \rightarrow i$ be an edge of level $r$. \newline
If $u^{\ast}_{a \rightarrow i} = 1$ the SAT formula defined by $\tau_{a \rightarrow i}$ is not satisfiable.

\begin{proof} Induction on $r$.
\begin{itemize}
	\item If $r = 0$, $a$ is a leaf and $\tau_{a \rightarrow i}$ defines a formula that consists of one empty clause.
	\item If $r = 1$ the message $u^{\ast}_{a \rightarrow i}$ cannot be $1$: $a$ contains one or more variables $j \neq i$ which all are leaf nodes that pass the message $m_{j \rightarrow a} = 0$ to $a$. The update equation yields $m_{a \rightarrow i} = \prod_j \theta(0) = 0$.
	\item Let $r \geq 2$ and $u^{\ast}_{a \rightarrow i} = -1$. For each variable $j \neq i$ of $a$ there is at least one $b$ with $J_j^b J_j^a u^\star_{b \rightarrow j}= 1$ or equivalently $J_j^b \neq J_j^a$ and  $u^\star_{b \rightarrow j} = 1$. By induction the subproblem $\tau_{b \rightarrow j}$ is not satisfiable and any assignment in which $j$ supports $a$ must violate at least one clause in $\tau_{b \rightarrow j}$.
\end{itemize}
\end{proof}
\end{lemma}



\subsection{Decimation Algorithm}

The algorithm that actually computes satisfying assignments based on the WP results is called \emph{Warning Inspired Decimation} (CID). CID uses WP as a subprocedure. For a satisfiable formula $\mathcal{F}$ after running WP, a partial assignment is computed that can be applied to $\mathcal{F}$ to obtain a smaller formula. Running CID recursively on that smaller formula gives a complete satisfying assignment. CID detects after the first run of WP if $\mathcal{F}$ is not satisfiable. \newline

To determine if a formula is satisfiable one has to compute certain values using the converged WP messages.

The \emph{local field} of a variable $i$ is defined as $H_i := - \sum_{b \in V(i)} J_i^b u^\ast_{b\rightarrow i}$. \newline $H_i$ is computed similar to $h_{i \rightarrow a}$ in the WP algorithm. While 
$h_{i \rightarrow a}$ gave the current tendency of $i$ when ignoring $a$, the local field of $i$ uses the converged messages and counts how often $i$ is fixed to $1$ or $0$ by all of its neighbours (including $a$). \newline
The \emph{contradiction number} $c_i$ of a variable is $1$ if $i$ gets fixed to different values and $0$, if not. $c_i$ can be obtained from the converged WP messages. If after running WP a variable has a contradiction number of $1$ the formula is unsatisfiable. \newline
The opposite direction is also true: If no contradiction appears, the formula is satisfiable and CID returns a valid assignment. 
After one run of WP a partial assignment can be obtained by setting all variables with active warnings to the value they were fixed to: $1$ if $H_i > 0$ and $0$ if $H_i < 0$. If $i$ received at least one warning, $H_i$ can not be $0$ because all incoming warnings fix $i$ to the same value. If no variable received a warning (all $H_i$ are $0$), CID chooses a random variable and sets it to a random value. \newline
Now the factor graph is modified (\emph{cleaned}) according to the partial assignment. The graph is cleaned by removing all fixed variables and all clauses that are satisfied by them.  The modified formula will be satisfiable and cause no contradictions in the following WP runs.
\newpage

\begin{lstlisting}[mathescape=true, frame = single]
	Warning Inspired Decimation
	  Input: A CNF formula and its factor graph 
	1. While there are unfixed variables
		1.1 Run WP
		1.2 If WP does not converge, return UN-CONVERGED. Else compute the numbers $H_i$ and $c_i$
		1.3 If at least one $c_i$ is $1$ return UNSAT.
		1.4 Choose a partial assignment and clean the graph
	2. Return the complete assignment
\end{lstlisting}

%\begin{lemma}
%If $u^{\star}_{a \rightarrow i} = 1$ the subproblem $a$ can not be satisfied if $i$ is removed from the set of its variables

%\begin{proof} 
%\begin{itemize}
%\item[] If the edge $(a, i)$ has level $1$, the clause $a$ only contains the variable $i$. Without $i$, it is an empty disjunction and evaluates to $0$
%\item[] If the edge $(a, i)$ has level $2$ the remaining variables $j \in V(a) \setminus i \neq \emptyset$ only appear in $a$. Their cavity fields must be $0$ so in any step of the algorithm $u^{\star}_{a \rightarrow i}$ is set to $0$.
%\end{itemize}
%\end{proof}
%\end{lemma}

\section{Belief Propagation}
\begin{frame}
	\frametitle{Belief Propagation Algorithm}
	\begin{itemize}
		\item General Algorithm for approximating marginal probabilities on factor graphs
		\item Applied in many different areas
		\item Can be viewed as generalization of Warning Propagation
	\end{itemize}
\end{frame}

\begin{frame}
	\frametitle{Belief Propagation Algorithm}
	Probability space built by all valid variable assignments $X = (x_1, \ldots, x_n)$ of a formula $\mathcal{F}$.
	\begin{itemize}
		\item Uniform distribution on all assignments 
		$$P(X) = \frac{1}{2^n}$$
		\item Restriction on valid assignments
	$$P(X \; | \; X \text{ satisfies } \mathcal{F}) = \frac{1}{\mathcal{N}} \prod_{a \in A} f_a(X)$$
	\item Goal: Compute $\mu_i := P(x_i = 1 \; X \text{ satisfies } \mathcal{F})$
	\end{itemize}
\end{frame}

\begin{frame}
	\frametitle{Messages}
	
	\begin{itemize}
		\item WP: $u_{a \rightarrow i}$ is $1$ if all neighbours $j \in V(a) \setminus i$ violate $a$
		\item BP: $\delta_{a \rightarrow i}$ is the \emph{probability} for this event
	\end{itemize}
	
	
\end{frame}

\begin{frame}
	\frametitle{Messages}
	
	\begin{itemize}
		\item WP: $h_{i \rightarrow a}$: preferred value of $i$
		\item BP: $\gamma_{i \rightarrow a}$: \emph{probability} for $i$ to violate $a$
	\end{itemize}
\end{frame}

\begin{frame}
	\frametitle{Decimation}
	Decimation Algorithm is pased on computed marginals
	\begin{itemize}
		\item Similiar to Warning Propagation
		\item Run BP Algorithm
		\item Assign $x_i = 1$ if $\mu_i > 0.5$ and $x_i = 0$ otherwise
		\item Clean the graph
	\end{itemize}
\end{frame}



\section{Conclusion}
\frame{
	\frametitle{Conclusion}
}

\end{document}